\subsection{Pass 10951: four clock ticks are the central spin sign}

The two negacyclic clock gauges of Pass 10948 have unique lifts in
$GL(2,3)$,
\[
 g_- = \begin{pmatrix}0&1\\1&1\end{pmatrix},
 \qquad
 g_+ = -g_-,
\]
with minimal polynomials $x^2-x-1$ and $x^2+x-1$.  Both have determinant
$-1$ and order eight, project to the same four-cycle of the clock rays, and
satisfy
\[
 g^2\in Q_8\subset SL(2,3),\qquad g^4=-I,\qquad g^8=I.
\]
An explicit all-products isomorphism identifies
$SL(2,3)$ with the binary tetrahedral group
$2T\subset\mathrm{Spin}(3)$ and sends $-I$ to quaternion $-1$.
Thus the determinant character is the finite Pin-component/Hesse-spinor-norm
bit, whereas the central $-I$ is the distinct $2\pi$ spin sign.

Under the standard inclusion $\mathrm{Spin}(3)\subset\mathrm{Spin}(9)$,
that central sign is the same one that Pass 10950 proves to be matter parity
on the Peirce $16$.  Hence four algebraic clock ticks reach the spin central
sign while eight return to the identity.  This identifies only the central
character; no full determinant-minus-one tick action on the Albert spinor,
continuum Pin structure, or physical time evolution is claimed.
